\documentclass[10pt,a4paper]{amsart}
\usepackage[margin=19mm]{geometry}
\usepackage{amsmath,amssymb,mathtools}
\usepackage[scr=rsfs,cal=euler]{mathalfa}
\usepackage{xfrac}
\usepackage[unicode,hidelinks,bookmarks=false]{hyperref}
\newcommand{\ie}{i.e.\ }
\newcommand{\cf}{cf.\ }
\newcommand{\resp}{respectively\ }
\newcommand{\Subsection}{\subsection}
\providecommand{\nobreakdash}{\nobreak\hbox{-}}
\setlength{\emergencystretch}{2em}
\multlinegap0pt
\usepackage{bm}
\usepackage{bbm}
\usepackage{dsfont}
\usepackage{xspace}
\usepackage{cancel}
\usepackage{mathtools}
\usepackage{amsthm}
\makeatletter
\@ifundefined{theorem}{\newtheorem{theorem}{Theorem}}{}
\@ifundefined{prop}{\newtheorem{prop}[theorem]{Proposition}}{}
\makeatother
\newcommand{\K}{\mathbb{K}}
\newcommand{\Z}{\mathbb{Z}}
\newcommand{\sL}{\mathcal{L}}
\title[Spectral sequences, Massey products and homology of covering]{Spectral sequences, Massey products and homology of covering spaces}
\author{Yongqiang Liu, Laurentiu Maxim, Botong Wang}
\date{Source: arXiv:2511.11893v1 (CC BY 4.0)}
\begin{document}
\maketitle

\noindent\emph{Source and licence.} Spectral sequences, Massey products and homology of covering spaces. Version arXiv:2511.11893v1 (CC BY 4.0), distributed under the Creative Commons Attribution 4.0 International licence, as recorded in the arXiv licence metadata for this exact version and shown on the abstract page https://arxiv.org/abs/2511.11893v1. Full rights notice: licenses/arxiv-2511.11893-CC-BY-4.0.txt.\par\smallskip

\noindent\textit{Editorial scope.} Counted unit: the Prop whose source label is \texttt{prop key} in the article (source0.tex lines 339--344, source label \texttt{prop key}), delivered together with the article's own complete original proof of it (source0.tex lines 345--403, one \texttt{proof} environment of 59 lines). No mathematics is written, repaired or re-derived: statement and proof are the article's own text line for line. Required context retained uncounted: cocycle (lines 329--329). Every definition, display and label of those lines is retained unchanged. Omitted from this package and not redistributed: 1--328, 330--338, 404--792. Nothing inside a retained excerpt is omitted or altered: every retained body is delivered exactly as the article's lines are written, so no omission receipt of any kind accompanies this package. Citations: the retained excerpts carry no citation command at all, so no bibliography entry of the article is transcribed and no reference is redistributed. Reference adaptation: none was needed and none was made; every reference inside the delivered unit resolves inside it, so no unresolved reference or citation is produced. Printed slips kept literal: no printed word, symbol, hypothesis or number of the counted statement or of its proof was altered. Layout and class: the article's own class is \texttt{amsart} with packages including geometry, amsmath, amscd, amssymb, amsfonts, latexsym, wasysym, mathrsfs, mathtools, hhline, color, xy, csquotes, url; the wrapper is the lane's pinned \texttt{amsart} editorial wrapper at 10pt on A4, which restates the theorem-like environments the retained text needs and adds the article's own macro definitions that the retained text needs (\texttt{\textbackslash K}, \texttt{\textbackslash Z}, \texttt{\textbackslash sL}), with extra wrapper packages (bm, bbm, dsfont, xspace, cancel, mathtools). The delivered numbering is the unit's own. Assets: no retained span contains a figure environment, a graphics-inclusion command, a PDF-page inclusion, a table or a TikZ picture; no image file is copied into this package, the package declares no asset dependency and no image file is redistributed with it. Licence: version arXiv:2511.11893v1 of this article is distributed under the Creative Commons Attribution 4.0 International licence, as recorded in the arXiv licence metadata for this exact version and shown on the abstract page https://arxiv.org/abs/2511.11893v1. A full rights notice is delivered at licenses/arxiv-2511.11893-CC-BY-4.0.txt. Characters: every retained excerpt is pure ASCII, so the delivered engine, XeLaTeX with its default Latin Modern text font, reports no missing character.
\begin{equation}\label{cocycle} \eta=f_\nu^*(\eta_0) \ \text{ and } \ \eta'=p^*(\eta) \end{equation} 

\begin{prop} \label{prop key}
Under the above notations, $(\partial+s\eta)^2=0$. In particular, $C_*(X, \eta, m)$ is a chain complex. Moreover, for any integer $i$ there is an isomorphism of $R/(s^{m+1})$-modules
\begin{equation}\label{eq1}
H_i\big(X, \sL^\nu\otimes_R R/(s^{m+1})\big)\cong H_i\big(C_*(X, \eta, m)\big).
\end{equation}
\end{prop}

\begin{proof}
We prove the first statement by explicit computation. For a simplicial chain $\omega$, we have
\begin{align*}
(\partial+s \eta)^2\omega&=\partial^2\omega+s\partial(\omega\cap \eta)+s \partial\omega\cap \eta+s^2\omega\cap(\eta\cup\eta)\\
&=\partial^2\omega-s\partial\omega\cap\eta+s\omega\cap\delta\eta+s\partial\omega\cap\eta+s^2\omega\cap(\eta\cup\eta)\\
&=\partial^2\omega+s\omega\cap\delta\eta+s^2\omega\cap(\eta\cup\eta)\\
&=0.
\end{align*}
Here we use the fact that $\eta$ is a $1$-cocycle on $X$, and also $\eta \cup \eta=0$ since $\eta$ is the pull-back of a $1$-cocycle on $S^1$ (cf. \eqref{cocycle}).

For the second part, we construct an isomorphism of complexes of $R/(s^{m+1})$-modules:
\begin{equation}\label{eq3}
\Phi: \left(C_*(X)\otimes_{\K}R/(s^{m+1}), \partial+s\eta \right)\to \left(C_*(X^\nu)\otimes_{R} R/(s^{m+1}), \partial_\nu\right),
\end{equation}
where $\partial$ and $\partial_\nu$ are the boundary maps of the simplicial chain complexes $C_*(X)$ and $C_*(X_\nu)$, respectively.

We first define $\Phi$ as a graded $R/(s^{m+1})$-module isomorphism, and then check its compatibility with the boundary maps. Let $\Delta\hookrightarrow X$ be a simplex of $X$. Since $f'_\nu$ is a map of simplicial complexes, there is a unique lift $\Delta'\hookrightarrow X^\nu$ such that $f'_\nu(\Delta'^\circ)\subset [0, 1)$, where $\Delta'^\circ$ denotes the interior of $\Delta'$. 
In particular, for a vertex $v$ of $X$, $v'\in X^\nu$ is the lifting such that $f'_\nu(v')\in [0, \frac{N-1}{N}]$.
Let $\Phi(\Delta\otimes 1)=\Delta'\otimes 1$. Notice that the set of $\Delta'$ defined above form a $R$-module basis of $C_*(X^\nu)$. Hence $\Phi$ is a graded $R/(s^{m+1})$-module isomorphism. 

Next, we prove that $\Phi$ is compatible with the boundary maps. We need to check that
\[
\partial_\nu\Phi(\Delta)=\Phi\big((\partial+s\eta)\Delta\big),
\]
where $\Delta$ is a $k$-simplex of $X$. 
If $f_\nu(\Delta)\neq [\frac{N-1}{N}, 1]$, then $\Delta\cap \eta=0$ and the above equality follows from the functoriality of chain complexes. 
So we only need to consider the case when $f_\nu(\Delta)=[\frac{N-1}{N}, 1]$. 

We choose the total ordering of vertices of $S^1$ by $x_0<x_1<\cdots <x_{N-1}$, and choose a total ordering of vertices of $X$, such that $f\colon X\to S^1$ preserves the ordering of the vertices. Furthermore, choose a total ordering of vertices of $X^\nu$ such that $p\colon X^\nu\to X$ preserves the ordering of the vertices. Write $\Delta=\langle v_0, \ldots, v_k\rangle$ according to the chosen ordering of vertices of $X$. 
 
Since $f_\nu$ preserves the order of the vertices of $X$ and $S^1$, and moreover $f_\nu(\Delta)=[\frac{N-1}{N}, 1]$, there exists $j\in \{0, \ldots, k-1\}$ such that ,
without any loss of generality, we may assume that
\[
f(v_0)=\cdots=f(v_j)=\frac{N-1}{N}\quad \text{and}\quad f(v_{j+1})=\cdots=f(v_k)=1.
\]
Consider the case when $j\geq 1$. In this case, 
\[
\Delta\cap \eta=(\langle v_0, v_1\rangle\cap \eta)\langle v_1, \dots, v_k\rangle=\left(\left\langle \frac{N-1}{N}, \frac{N-1}{N}\right\rangle\cap \eta_0\right)\langle v_1, \dots, v_k\rangle=0.
\]
Denote the action of $1\in \Z$ on $X^{\nu}$ by $\sigma$. Then
\begin{align*}
&\Phi\big((\partial+s\eta)\Delta\big)=\Phi(\partial\Delta)\\
=&\sum_{0\leq i\leq j}(-1)^i\Phi\big(\langle v_0, \ldots, \widehat{v_i}, \ldots, v_j, \ldots v_k\rangle\big)+\sum_{j+1\leq i\leq k}(-1)^i\Phi\big(\langle v_0, \ldots, v_j, \ldots, \widehat{v_i}, \ldots, v_k\rangle\big)\\
=&\sum_{0\leq i\leq j}(-1)^i\left\langle v'_0, \ldots, \widehat{v'_i}, \ldots, v'_j, \sigma(v'_{j+1}), \ldots \sigma(v'_k)\right\rangle\\
&+\sum_{j+1\leq i\leq k}(-1)^i\left\langle v'_0, \ldots, v'_j, \sigma(v'_{j+1}), \ldots, \widehat{\sigma(v'_i)}, \ldots, \sigma(v'_k)\right\rangle\\
=&\partial_\nu\langle v'_0, \ldots, v'_j, \sigma(v'_{j+1}), \ldots \sigma(v'_k)\rangle\\
=&\partial_\nu\Phi(\Delta).
\end{align*}

Consider the case when $j=0$. In this case, $\Delta\cap \eta=\langle v_1, \ldots, v_k\rangle$. Hence,
\begin{align*}
\Phi((\partial+s\eta)\Delta)=&\sum_{1\leq i\leq k}(-1)^i\Phi\big(\langle v_0, v_1, \ldots, \widehat{v_i}, \ldots v_k\rangle\big)+\Phi\big(\langle v_1, \ldots, v_k\rangle\big)+s\Phi\big(\langle v_1, \ldots, v_k\rangle\big)\\
=&\sum_{1\leq i\leq k}(-1)^i\Phi\big(\langle v_0, v_1, \ldots, \widehat{v_i}, \ldots v_k\rangle\big)+t\Phi\big(\langle v_1, \ldots, v_k\rangle\big)\\
=&\sum_{1\leq i\leq k}(-1)^i\left(\Big\langle v'_0, \sigma(v'_1), \ldots, \widehat{\sigma(v'_i)}, \ldots \sigma(v'_k)\Big\rangle\right)+t\big(\langle v'_1, \ldots, v'_k\rangle\big)\\
=&\sum_{1\leq i\leq k}(-1)^i\left(\left\langle v'_0, \sigma(v'_1), \ldots, \widehat{\sigma(v'_i)}, \ldots \sigma(v'_k)\right\rangle\right)+\big(\big\langle \sigma(v'_1), \ldots, \sigma(v'_k)\big\rangle\big)\\
=&\partial_\nu \big\langle v'_0, \sigma(v'_1), \ldots,  \sigma(v'_k)\big\rangle\\
=&\partial_\nu\Phi(\Delta).\qedhere
\end{align*}
\end{proof}

\end{document}
